%=====================================================================
\section{Simulation Test}
%=====================================================================
\textit{Owner Margarita Kholostova. Checked by Aarav Jain.}


\textbf{What we test.} We run our own ALS (\texttt{my\_als} in \texttt{src/methods.py}, which solves every small system with our own Gaussian elimination) on the Iteration 2 simulation. The data, $k = 3$, $\lambda = 1$, 50 steps and seed 0 are the same as in Iteration 2. We compare the result with two things, the true matrix and a library.

\subsection{Compare with the truth}

The truth is $R_{\text{true}}$ on the 5{,}719 hidden cells. Our ALS has RMSE 0.147 on these cells, which is below the pass line 0.3. Most of this error comes from the penalty and not from the noise: $\lambda = 1$ pulls every $u_u$ and $v_i$ a little toward 0, and a user here has only about 23 ratings to push back. With $\lambda = 0.01$ the same run has RMSE 0.059 (\texttt{src/check\_als\_steps.py}, part 3). So the method finds the hidden values, and $\lambda = 1$ gives up a little accuracy here for protection against overfitting, which matters on the real data (Question 4).

\subsection{Compare with a library}

The library version is the Iteration 2 function \texttt{als}, which solves the same systems with \texttt{numpy.linalg.solve}. Both start from the same random $U$ and $V$. After 50 steps the largest difference between the two matrices $UV^\top$ is $7.1\times10^{-15}$. For one user the two solutions of $(V_u^\top V_u + \lambda I)u_u = V_u^\top r_u$ differ by less than $10^{-16}$, and putting our $u_u$ back into the system gives a largest residual of $1.8\times10^{-15}$. The last digits of these tiny numbers change from one computer to another; Question 2 repeats the check for every user and movie. These numbers are rounding error: both codes compute the same formula from the same data, and a float keeps about 16 digits. So our Gaussian elimination is correct.

\subsection{Error against the size of the matrix}

We repeat the simulation for four sizes with the same shape and the same 30\% of known cells (Table~\ref{tab:sizes}, Figure~\ref{fig:sizes}).

\begin{table}[h]
\centering
\begin{tabular}{lll}
\hline
Size $n_u \times n_i$ & Known cells & RMSE on hidden cells \\
\hline
$50 \times 40$ & 592 & 0.414 \\
$100 \times 80$ & 2{,}281 & 0.147 \\
$200 \times 160$ & 9{,}601 & 0.057 \\
$400 \times 320$ & 38{,}283 & 0.033 \\
\hline
\end{tabular}
\caption{Error against the truth for bigger matrices (from \texttt{src/check\_simulation.py}).}
\label{tab:sizes}
\end{table}

\begin{figure}[h]
  \centering
  \includegraphics[width=0.55\textwidth]{simulation_error_vs_size.png}
  \caption{More data gives better guesses. Error of ALS on the hidden cells for four matrix sizes, with 30\% of cells known in each.}
  \label{fig:sizes}
\end{figure}

\textbf{What Figure~\ref{fig:sizes} shows.} Each dot is one run of our ALS on a simulated matrix where we know the truth. The horizontal axis is the number of users (50, 100, 200, 400, with 80 movies for every 100 users). The vertical axis is the RMSE on the hidden cells, so lower is better. Both axes are on a log scale, so a straight line means the error follows a power of the size. The dashed line is the noise we added (0.1).

\textbf{What we see.} The error falls from 0.414 to 0.147, 0.057 and 0.033. From 100 users up the dots lie almost on a straight line with slope $-1.08$. The smallest matrix is above the pass line 0.3, because a user there has only about 12 known ratings for 3 unknown numbers.

\textbf{Why.} If the noise were the only source of error, the error would be about $\sigma\sqrt{k(n_u+n_i)/N}$, the noise times the square root of the number of unknowns over the number of known cells. When $n_u$ doubles, $n_u+n_i$ doubles and $N$ grows 4 times, so this falls with slope $-1/2$ (it is 0.049, 0.034 and 0.024 for the three larger sizes). Our slope is steeper because of the penalty: $\lambda = 1$ pulls each $u_u$ toward 0 by about $\lambda/|\Omega_u|$ of its size, and $|\Omega_u|$ grows like $n_u$, so this part of the error falls with slope $-1$. With a small penalty, $\lambda = 0.01$, the errors are 0.059, 0.036 and 0.025 and the slope is $-0.63$, close to the noise rate (\texttt{src/check\_als\_steps.py}, part 3).

\textbf{Conclusion.} The error shrinks as the data grows, and both parts of it shrink at the rate the math predicts: the pull of the penalty like $1/n_u$ and the noise like $1/\sqrt{n_u}$. So our ALS is a correct estimator and not just a good fit to one matrix. The error even goes below the noise level 0.1, because it is measured against $R_{\text{true}}$, which has no noise.

\subsection{Stress test with fewer known cells}

Our method assumes that every user and every movie has enough known ratings. We break this on purpose and lower the share of known cells (Table~\ref{tab:stress}, Figure~\ref{fig:stress}).

\begin{table}[h]
\centering
\begin{tabular}{llll}
\hline
Share known & Ratings per user (mean, min) & RMSE on hidden cells & $J$ never increases \\
\hline
0.30 & 22.8, 10 & 0.147 & yes \\
0.20 & 15.3, 6 & 0.263 & yes \\
0.15 & 11.5, 4 & 0.416 & yes \\
0.10 & 8.2, 1 & 1.372 & yes \\
0.07 & 6.0, 1 & 1.857 & yes \\
0.05 & 4.2, 0 & 1.810 & yes \\
\hline
\end{tabular}
\caption{Stress test on the $100 \times 80$ simulation (from \texttt{src/check\_simulation.py}).}
\label{tab:stress}
\end{table}

\begin{figure}[h]
  \centering
  \includegraphics[width=0.55\textwidth]{stress_test_known_share.png}
  \caption{With too few known ratings ALS stops working. Error of ALS on the hidden cells of the $100 \times 80$ simulation as the share of known cells goes down.}
  \label{fig:stress}
\end{figure}

\textbf{What Figure~\ref{fig:stress} shows.} Each dot is one run of our ALS on the same simulated matrix, with a different share of cells kept (30\% down to 5\%). The horizontal axis is the share of known cells, and the vertical axis is the RMSE on the hidden cells. The dashed line is the pass line 0.3 from the Iteration 3 instructions, so dots below it pass.

\textbf{What we see.} ALS stops recovering the matrix between 20\% and 15\% of known cells: at 20\% it still passes (0.263) and at 15\% it does not (0.416). Between 15\% and 10\% the error jumps to 1.372, as bad as the filled SVD (1.394), and then stays high.

\textbf{Why.} The reason is in the system for one user. The vector $u_u$ has $k = 3$ unknown numbers, so it needs at least 3 ratings to be fixed by the data. At 10\% one user has a single rating and at 5\% one user has none, and then $V_u^\top V_u$ has rank below 3 and only $\lambda I$ keeps the system solvable (Iteration 2, Question 3). The system still has a solution, but the solution comes mostly from the penalty and not from the data.

In all six runs $J$ never increases. So Question 4 of Iteration 2 holds even when the method fails: a small $J$ on the known cells does not mean good guesses on the hidden cells.

\textbf{Conclusion.} ALS needs enough ratings per user and per movie, at least about $k$ of them and better several times $k$. On this matrix that means at least 20\% of cells known. MovieLens has only 6.3\%, but its matrix is much larger, so a typical user still has many ratings (median 65); the users and movies with very few ratings are where we expect the worst guesses.

%=====================================================================
\section{Real Data}
%=====================================================================
\textit{Owner Margarita Kholostova. Checked by Aarav Jain.}


\subsection{Cleaning and split}

We use the split that MovieLens provides, \texttt{u1.base} for training and \texttt{u1.test} for testing, as the instructions ask. Table~\ref{tab:clean} lists every step.

\begin{table}[h]
\centering
\begin{tabular}{ll}
\hline
Step & Number \\
\hline
Ratings in \texttt{u1.base} (train) & 80{,}000 \\
Ratings in \texttt{u1.test} (test) & 20{,}000 \\
Rows removed & 0 \\
Cells that are in both train and test & 0 \\
Size of both matrices & $943 \times 1682$ \\
Transformations of the ratings & none (stars 1 to 5 as given) \\
Users with no train rating & 0 \\
Fewest train ratings for one user & 4 \\
Movies with no train rating & 32 \\
Test ratings on these 32 movies & 32 \\
\hline
\end{tabular}
\caption{Cleaning steps and the split (from \texttt{src/run\_real.py}).}
\label{tab:clean}
\end{table}

For the 32 movies with no train rating the system for $v_i$ is $\lambda I\, v_i = 0$, so $v_i = 0$ and the guess is 0 stars. We keep these 32 test ratings, because removing them would make our error look better than it is. They are 0.16\% of the test set.

%=====================================================================
\section{Iteration 3 Math Questions}
%=====================================================================

\subsection*{Question 3 Truncated SVD baseline with the mean}
\textit{Owner Margarita Kholostova. Checked by Aarav Jain.}

\textbf{Step 1. Fill.} Let $\mu$ be the mean of the known train ratings ($\mu = 3.528$). We build a full matrix $\tilde R$ with $\tilde r_{ui} = r_{ui}$ if $(u,i) \in \Omega$ and $\tilde r_{ui} = \mu$ otherwise.

\textbf{Step 2. SVD.} $\tilde R$ has no empty cells, so it has an SVD $\tilde R = P \Sigma Q^\top$.

\textbf{Step 3. Cut.} We keep the $k$ largest singular values, $\tilde R_k = P_k \Sigma_k Q_k^\top$. By the Eckart--Young theorem (Iteration 2, Question 5) $\tilde R_k$ is the closest rank-$k$ matrix to $\tilde R$ in the Frobenius norm. The guess for user $u$ and movie $i$ is the entry $(u,i)$ of $\tilde R_k$.

\textbf{Result.} On the simulation the filled SVD has RMSE 1.394 on the hidden cells and ALS has 0.147. On MovieLens the test RMSE is 1.099, 1.085, 1.080 and 1.083 for $k = 2, 5, 10, 20$.

\textbf{Conclusion.} $\tilde R_k$ is the best fit to $\tilde R$, but 95\% of the cells of $\tilde R$ are the number $\mu$ that we put there ourselves (80{,}000 real ratings out of 1{,}586{,}126 cells). So the SVD mostly learns our filling and not the ratings, and on MovieLens it is even worse than the mean of each movie (1.033). The SVD here is \texttt{numpy.linalg.svd}, because it is only a baseline to compare with and not our method.

\subsection*{Question 4 Train and test split and the test RMSE}
\textit{Owner Margarita Kholostova. Checked by Aarav Jain.}

\textbf{Definition.} Let $\Omega_{\text{test}}$ be the set of the 20{,}000 test pairs. The method never sees them. We find $U$ and $V$ from the train pairs only and then compute
\[
\text{RMSE}_{\text{test}} = \sqrt{\frac{1}{|\Omega_{\text{test}}|} \sum_{(u,i) \in \Omega_{\text{test}}} \left(r_{ui} - u_u^\top v_i\right)^2 } .
\]
It is measured in stars, so an RMSE of 1 means that a typical guess is wrong by about one star.

\textbf{Baselines.} Guessing $\mu = 3.528$ for everything gives 1.154. Guessing the mean train rating of each movie gives 1.033 (a movie with no train rating gets $\mu$).

\textbf{Result.} Table~\ref{tab:grid} gives the test RMSE of our ALS after 50 steps with seed 0 (Figure~\ref{fig:grid}).

\begin{table}[h]
\centering
\begin{tabular}{llll}
\hline
 & $\lambda = 0.1$ & $\lambda = 1$ & $\lambda = 10$ \\
\hline
$k = 2$ & 0.950 & \textbf{0.940} & 0.991 \\
$k = 5$ & 1.071 & 0.973 & 0.984 \\
$k = 10$ & 1.332 & 1.074 & 0.986 \\
$k = 20$ & 1.539 & 1.201 & 0.998 \\
\hline
\end{tabular}
\caption{Test RMSE of ALS on \texttt{u1.test}. The movie mean baseline is 1.033 (from \texttt{src/run\_real.py}).}
\label{tab:grid}
\end{table}

\begin{figure}[h]
  \centering
  \includegraphics[width=0.6\textwidth]{real_test_rmse.png}
  \caption{A small $k$ or a large $\lambda$ works best on real data. Test RMSE of ALS on \texttt{u1.test} for each $k$ and $\lambda$. The dashed line is the movie mean baseline (1.033).}
  \label{fig:grid}
\end{figure}

\textbf{What Figure~\ref{fig:grid} shows.} It draws Table~\ref{tab:grid}. The horizontal axis is $k$, the number of factors for each user and movie. The vertical axis is the test RMSE in stars, so lower is better. Each line is one value of the penalty $\lambda$. A point below the dashed line beats the simple guess "the average rating of this movie".

\textbf{What we see.}
\begin{itemize}
  \item The best result is 0.940 at $k = 2$ and $\lambda = 1$. It is 9\% better than the movie mean (1.033), and 7 of the 12 settings beat this baseline.
  \item With a small $\lambda$ a bigger $k$ makes the test error worse. At $k = 20$ and $\lambda = 0.1$ the train RMSE is 0.470 but the test RMSE is 1.539. This is the overfitting from Step 3 of Question 1 in Iteration 2: some users have only 4 train ratings, and 20 free numbers fit 4 ratings perfectly.
  \item With $\lambda = 10$ the test error stays between 0.984 and 0.998 for every $k$. So the penalty does the job we added it for.
  \item We picked $k = 2$ and $\lambda = 1$ because they have the lowest test RMSE, so 0.940 is slightly optimistic: the test set was used to choose the setting. A fully fair number would choose $k$ and $\lambda$ on a validation part of \texttt{u1.base} and use \texttt{u1.test} only once. All seven settings that beat the movie mean are between 0.940 and 0.998, so the conclusion that ALS beats both baselines and the filled SVD does not depend on this choice.
\end{itemize}

\begin{sloppypar}\textbf{Check with a library.} At the best setting we ran the Iteration 2 \texttt{als} (with \texttt{numpy.linalg.solve}) from the same start. The largest difference in $UV^\top$ over all $943 \times 1682$ cells is $1.3\times10^{-13}$, and the two test RMSE values are the same to 10 digits (0.9398201282). The difference is larger than on the simulation ($7.1\times10^{-15}$), because the real matrix has 2{,}625 systems per step instead of 180 and the ratings are larger numbers, so more rounding errors add up. It is still rounding error and not a sign or scaling convention: both codes start from the same $U$ and $V$, and each system has one solution (Iteration 2, Question 3), so there is nothing to align.\end{sloppypar}

\textbf{Conclusion of Figure~\ref{fig:grid} and Table~\ref{tab:grid}.} More factors are not always better: a bigger $k$ helps only if $\lambda$ grows with it, because with few ratings per user extra factors fit the noise. On ratings that the method never saw, our ALS guesses within about 0.94 stars. This is better than both simple baselines and than the filled SVD. So the method works on real data and not only on the simulation, if $k$ is small or $\lambda$ is large enough.

%=====================================================================
\section{Verification Table}
%=====================================================================
\textit{Owner Margarita Kholostova. Checked by Aarav Jain.}

% rows for the simulation and the real data

\begin{table}[h]
\centering
\small
\resizebox{\textwidth}{!}{%
\begin{tabular}{lllll}
\hline
Data & Quantity & Compared with & Difference & Pass? \\
\hline
Simulation & $UV^\top$, hidden cells & $R_{\text{true}}$ & RMSE $= 0.147$ & yes \\
Small example & $U^{(1)}, V^{(1)}$ from \texttt{methods.py} & exact fractions & $\max|\cdot| = 0$ & yes \\
Small example & $J$ after the movie update & $J$ after the user update (39) & $212/7 \approx 30.29$ & yes \\
Simulation & $UV^\top$ & ALS with numpy solve & $\max|\cdot| = 7.1\times10^{-15}$ & yes \\
Simulation & all 180 systems of one step & numpy solve & $\max|\cdot| = 8.9\times10^{-16}$ & yes \\
Simulation & $J$ after each of 100 half steps & never increases & largest change $-1.1\times10^{-5}$ & yes \\
Simulation & largest step of $J$ & $\le 0$ & $-2.4\times10^{-5}$ & yes \\
Simulation & filled SVD, hidden cells & ALS (0.147) & RMSE $= 1.394$ & yes \\
Stress test (10\% known) & $UV^\top$, hidden cells & $R_{\text{true}}$ & RMSE $= 1.372$ & fails \\
Real data & $UV^\top$ & ALS with numpy solve & $\max|\cdot| = 1.3\times10^{-13}$ & yes \\
Real data & all 2{,}625 systems of one step & numpy solve & $\max|\cdot| = 5.0\times10^{-15}$ & yes \\
Real data & $J$ after each of 100 half steps & never increases & largest change $-17.4$ & yes \\
Real data & test RMSE & movie mean (1.033) & 0.940 & yes \\
Real data (12 runs) & largest step of $J$ & $\le 10^{-9}$ & 12 of 12 & yes \\
\hline
\end{tabular}}
\caption{Verification table for the simulation and the real data.}
\label{tab:verify}
\end{table}
